%% LyX 2.4.4 created this file.  For more info, see https://www.lyx.org/.
%% Do not edit unless you really know what you are doing.
\documentclass[english]{article}
\usepackage[T1]{fontenc}
\usepackage[utf8]{inputenc}
\setcounter{secnumdepth}{-2}
\setlength{\parindent}{0bp}
\usepackage{amsmath}
\usepackage{amsthm}
\usepackage{amssymb}

\makeatletter
%%%%%%%%%%%%%%%%%%%%%%%%%%%%%% Textclass specific LaTeX commands.
\theoremstyle{plain}
\newtheorem*{lem*}{\protect\lemmaname}

\@ifundefined{date}{}{\date{}}
%%%%%%%%%%%%%%%%%%%%%%%%%%%%%% User specified LaTeX commands.
\usepackage{graphicx}
\usepackage{geometry}
\newcommand{\limi}{\underset{n\rightarrow\infty}{\lim}}
\newcommand{\limxz}{\underset{x\rightarrow x_{0}}{\lim}}
\newcommand{\limfxz}{\underset{x\rightarrow x_{0}}{\lim}f(x)}
\newcommand{\limfxm}{\underset{x\rightarrow x_{0}^{-}}{\lim}f(x)}
\newcommand{\limfxp}{\underset{x\rightarrow x_{0}^{+}}{\lim}f(x)}
\newcommand{\limxm}{\underset{x\rightarrow x_{0}^{-}}{\lim}}
\newcommand{\limxp}{\underset{x\rightarrow x_{0}^{+}}{\lim}}
\newcommand{\sqrm}{M_{m\times m}(\mathbb{F})}
\newcommand{\seqa}{(a_{n})_{n=1}^{\infty}}
\newcommand{\seqx}{(x_{n})_{n=1}^{\infty}}
\newcommand{\funcdr}{f:D\rightarrow\mathbb{R}}
\newcommand{\der}{\limxz\frac{f(x)-f(x_{0})}{x-x_{0}}}
\usepackage[all]{pdfprivacy}

\makeatother

\usepackage{babel}
\providecommand{\lemmaname}{Lemma}

\begin{document}
\title{{\Large Linear Algebra Done Right by Axler, Result 1.45\vspace{-2em}}}
\maketitle
\begin{lem*}
condition for a direct sum

Suppose $V_{1},\dots,V_{m}$ are sub-spaces of $V$. Then $V_{1}+\dots+V_{m}$
is a direct sum if and only if the only way to write $0_{V}$ as a
sum $v_{1}+\dots+v_{m}$ s.t $\forall1\leq k\le m,\;v_{k}\in V_{k}$
is with $\forall1\leq k\le m,\;v_{k}=0_{V}$.
\end{lem*}
\begin{proof}
We prove the two implications separately.

\medskip{}

\textbf{First direction:}

Suppose $V_{1}+\dots+V_{m}$ is a direct sum.

Then, in particular, $0\in V_{1}+\dots+V_{m}$ can be written in only
one way as a sum $v_{1}+\dots+v_{m}=0$ where $\forall1\leq k\leq m,\ v_{k}\in V_{k}$.

Since $V_{1},\dots,V_{m}$ are all sub-spaces of $V$, they all contain
$0$ (i.e.\ $0_{V}$).

And since $0+\dots+0=0$, we get that this is the one way as needed.

\medskip{}

\textbf{Second direction:}

Let us assume $0+\dots+0=0$ is the only way to write $0$ as a sum
$v_{1}+\dots+v_{m}$ s.t.\ $\forall1\leq k\leq m,\ v_{k}\in V_{k}$.

Then for every $a\in V_{1}+\dots+V_{m}$, if $v_{1}+\dots+v_{m}=a$
s.t.\ $\forall1\leq k\leq m,\ v_{k}\in V_{k}$ and $v'_{1}+\dots+v'_{m}=a$
s.t.\ $\forall1\leq k\leq m,\ v'_{k}\in V_{k}$, we get that 
\[
v_{1}+\dots+v_{m}-\left(v'_{1}+\dots+v'_{m}\right)=a-a=0
\]
\[
\therefore\quad(v_{1}-v'_{1})+(v_{2}-v'_{2})+\dots+(v_{m}-v'_{m})=0
\]

each $V_{k}$ is a sub-space, hence closed under addition and scalar
multiplication, therefore for each $1\le k\le m$,

$(v_{k}-v'_{k})=v_{k}+(-1)\cdot v'_{k}\in V_{k}$, so the assumption
dictates that
\[
v_{1}-v'_{1}=0\ \land\ v_{2}-v'_{2}=0\ \land\ \dots\ \land\ v_{m}-v'_{m}=0
\]
\[
\therefore\quad v_{1}=v'_{1}\ \land\ v_{2}=v'_{2}\ \land\ \dots\ \land\ v_{m}=v'_{m}
\]
so there is exactly one way to write $a$ in this manner ($a\in V_{1}+\dots+V_{m}$
definitionally means one exists) and $V_{1}+\dots+V_{m}$ is a direct
sum as needed, hence the notation $V_{1}\oplus\dots\oplus V_{m}$
can be used.
\end{proof}

\end{document}
